\begin{resumo}[Abstract]

This work characterizes the probabilistic calibration of a Temporal Fusion Transformer (TFT) applied to multi-horizon daily return forecasting, compares it with a hierarchy of reference models, and describes the relative contribution of indicator families. The research problem stems from the fact that, for daily returns of liquid stocks, the predictable component of the mean is minimal, so evaluations restricted to point metrics tend to measure noise; the scientific object is therefore the predictive distribution rather than point accuracy. The scope is single-asset (one model per asset), with AAPL as the pilot asset and main horizons $h+1$ and $h+7$, with $h+30$ as supplementary. The methodology integrates market data, technical indicators, news sentiment estimated with FinBERT, and accounting fundamentals, with temporal causality verified feature by feature, and trains the TFT in quantile mode over a dense grid of levels, with monotone rearrangement of the quantiles. The experimental protocol uses expanding-window walk-forward validation with purging and embargo in trading days and a dedicated calibration set, multiple seeds and folds, strict alignment by target timestamp, and comparison against a pre-declared hierarchy of baselines: zero return, historical mean, AR(1), EWMA volatility, historical empirical quantiles, and quantile gradient boosting. Pinball loss is the primary metric and CRPS is complementary; inference relies on the Diebold-Mariano test (HAC/HLN, one-sided) with Holm adjustment and the Model Confidence Set, and native calibration is contrasted with conformalized quantile regression as an empirical-coverage benchmark. Feature-family contribution combines the Variable Selection Network, permutation importance, and ablation. The verdict is issued by a pre-registered, hashed scorecard, and refutation of the hypotheses is treated as a valid result. Confirmatory results will be consolidated after the confirmatory round is executed.

\textbf{[Version note — to be removed in the final defense version]} This abstract is a living version, updated at each advising round as confirmatory results are consolidated. The final version submitted to the examining committee will not include this note.

\textbf{Keywords}: Temporal Fusion Transformer. Quantile forecasting. Probabilistic calibration. Conformal prediction. Out-of-sample evaluation.

\end{resumo}
